\documentclass[10pt]{article}
\usepackage[margin=20mm]{geometry}
\usepackage{amsmath,amssymb,mathtools,fontspec,unicode-math,xurl,hyperref}
\IfFontExistsTF{DejaVu Serif}{\setmainfont{DejaVu Serif}}{\setmainfont{Latin Modern Roman}}
\setmathfont{latinmodern-math.otf}
\setlength{\parindent}{0pt}\setlength{\parskip}{6pt}\sloppy
\title{Initial scientific review: M25-0095
Compatibility of canonical \(\ell\)-adic local systems
on adjoint Shimura varieties}
\author{Reviewing assistant}
\date{26 October 2026}
\begin{document}
\maketitle
Abstract: The main compatibility argument was traced through the abstract extension/companion argument, superrigidity, construction of special points, Frobenius valuation, elimination of outer automorphisms, and rationality. No counterexample satisfying all the hypotheses of a main theorem was established. Two proof details need explicit clarification: the passage from non-isomorphic tori to non-isomorphic rational character representations, and the use of a maximal compact subgroup assertion for a torus relative to a fixed hyperspecial subgroup. Concrete algebraic checks explain each concern, and local remedies are given. Their unresolved scope is recorded as Needs Evidence, without a confirmed Major or Minor finding. This is initial review, not scientific acceptance.\par
\section*{Version, scope and main claims}
The bound article is Christian Klevdal and Stefan Patrikis, Inventiones mathematicae 242 (2025), 305–335, DOI 10.1007/s00222-025-01357-6, publisher Version of Record, published online 25 August 2025. The original PDF has 31 pages and 1,462,931 bytes. All references below use PDF page numbers; printed page equals PDF page plus 304. All 31 pages were read as full extracted text and actually viewed as page images. Reading the references is not verification of all cited works.\par
The source manifest contains the original PDF and publisher HTML, with no author TeX or released code identified. Both manifest files were hash-checked. This absence is not a defect. No author code was executed. The review’s calculations are small exact checks written independently by the reviewer.\par
Theorems 1.1 and 1.3 (pp. 3, 5), and the more general Theorem 3.10 and Corollary 3.12 (pp. 25–28), require a neat level, discreteness of \(Z_G(\mathbb Q)\) in \(Z_G(\mathbb A_f)\), and real rank at least two for every \(\mathbb Q\)-simple adjoint factor. The uniform exceptional integer is chosen before \(\ell,\ell'\) and the closed point. After specializing to a number-field point, a further integer may depend on that point. The conclusion concerns semisimple geometric conjugacy classes in the adjoint GIT quotient. It does not provide a rational representative in \(G^\mathrm{ad}(\mathbb Q)\), compatibility in the full central extension \(G\), or the missing motives. The qualifications in Remarks 1.2, 1.4, 1.5 and 3.14 were preserved. Results of subsequent papers cited there were not used to certify the present proof.\par
\section*{Checks of the proof chain}
\subsection*{Adjointness and descent}
Lemma 2.1 (p. 8) is correct by uniqueness of conjugators. Zariski density gives \[C_H(\rho(\Pi'))=Z(H)=1.\] For \(x,y\in\Pi\), the element \(h_xh_y\) conjugates \(\rho\) to \(\rho^{xy}\), so uniqueness gives \(h_{xy}=h_xh_y\). For \(x\in\Pi'\) it gives \(h_x=\rho(x)\). Applying every element of \(\operatorname{Gal}(K^s/K)\) gives \(h_x\in H(K)\). This also identifies exactly why the extension argument cannot simply be copied to a group with center: central characters of \(\Pi/\Pi'\) can produce distinct extensions.\par
In Proposition 2.3 (pp. 9–11), the hypothesis \(F(\lambda)=F\) is essential and is not proved for arbitrary rigid systems. Canonical models supply a common field in the Shimura application (p. 18). The use of \(N'=\gcd_\lambda N_\lambda\) is consistent: a prime outside this gcd is good for at least one coefficient place, and that place cannot lie over the same rational prime because its own coefficient prime was included in \(N_\lambda\).\par
The companion is pulled back by tame specialization, not by an asserted inverse to specialization. Its geometric Zariski closure is a normal subgroup \(H\) of the adjoint semisimple group. The quotient \(G/H\) is densely generated by a single Frobenius and is therefore commutative. A quotient of the adjoint semisimple group that is commutative is trivial. This verifies the density step needed for superrigidity. After geometric factorization, inertia centralizes the dense geometric image; adjointness makes inertia trivial. The final extension across the fiber still depends on the cited ramification detection and Hilbert irreducibility results. Their full proofs were not reproduced.\par
\subsection*{Removing the outer ambiguity}
Proposition 2.8 (pp. 12–13) correctly needs a Frobenius element fixed by no automorphism with nontrivial outer class. Equality of the semisimple classes says that a lift \(\widetilde\tau\) of the outer class carries a representative \(t\) to a conjugate. Choosing \(g\) with \((\operatorname{inn}(g)\widetilde\tau)(t)=t\) and transporting to the coefficient field produces an actual automorphism fixing \(t_\lambda\). Assumption 2.7(2) then forces the outer class to vanish. Merely knowing equality of characteristic polynomials in one representation would not suffice for this step.\par
Proposition 3.3 (pp. 16–17) reduces the dense representation to a product of arithmetic subgroups. The finite index of that product ensures density is retained. Commuting normal images in a simple adjoint target leave one nontrivial source factor per target factor. The resulting surjective algebraic endomorphism has finite central kernel, hence is an automorphism because the source is adjoint. Its agreement on the normal finite-index subgroup extends to the whole group by the trivial centralizer argument. For the final descent one may first enlarge the auxiliary field to an algebraically closed overfield; then invariance under automorphisms over the algebraically closed base gives descent. Alternatively, the unique conjugator is the unique geometric point of the conjugacy equations, which are defined over the base. This supplies the routine field-descent detail without assuming that an arbitrary field extension has a large automorphism group. Margulis’s external theorem itself remains an input.\par
\subsection*{Canonical systems and the special-point formula}
The right action, opposite automorphism group and base-point choices in pp. 14–18 were checked together. The topological inclusion is sent to \(a^{-1}\gamma a\), as shown by \([\gamma^{-1}x,a]_K=[x,\gamma a]_K\). The convention of geometric Frobenius agrees with the reciprocity convention on p. 8. Lemma 3.2 gives \[r_x(\alpha)=q\,a k_0(\sigma)a^{-1},\qquad qx=x.\] When the centralizer is the torus, the component away from the support of the uniformizer yields Equation (6), \[k_{0,\ell}(\operatorname{Frob}_w)=a_\ell^{-1}q^{-1}a_\ell.\] The inverse is necessary. The same rational \(q\) works for the coefficient primes, which is the required horizontal compatibility, subject to the torus and integrality inputs.\par
In Lemma 3.6 and Proposition 3.7 (pp. 19–22), the full Weyl-group action and transitive action on absolutely simple factors give irreducibility of the rational character representation. The local prescriptions at \(p\), at infinity and at finitely many pairwise auxiliary primes are finite in number. The twisting argument preserves the local torus isomorphism types outside one auxiliary prime. Its use of Brauer–Nesbitt and Chebotarev concerns rational character representations, not integral lattices. The Hasse principle, existence of the cohomology class and weak approximation for the variety of tori remain external inputs.\par
\subsection*{Frobenius closure and rationality}
For a rational element \(t\) of a \(\mathbb Q\)-irreducible torus, infinite order implies that the identity component of the \(\mathbb Q\)-Zariski closure of \(\langle t\rangle\) is a nonzero subtorus, hence the whole torus. With pairwise non-isogenous irreducible factors, a connected subtorus projecting onto all factors is their full product: its character quotient has no diagonal identification between distinct irreducible rational modules. This is the precise input required on p. 27; the distinction in the next section matters.\par
An automorphism fixing a maximal torus pointwise fixes its characters and roots. Over an algebraically closed field it acts on the one-dimensional root spaces by compatible scalars, and is conjugation by an element of that torus. Thus the final elimination of outer automorphisms is valid once the full-torus closure has been obtained. The base-change assertion in footnote 9 is appropriate because the powers of \(t\) are rational.\par
Finally, compatibility over all coefficient embeddings is used before descent from \(\overline{\mathbb Q}\) to \(\mathbb Q\). For a finite field of definition, invariance under the decomposition groups for almost all rational primes forces invariance under the entire finite Galois group by Chebotarev. This proves rationality in the quotient, not rational conjugacy in \(G^\mathrm{ad}\). Corollary 3.12 then follows because an \(F\)-rational point determines a Galois-stable geometric component, whose field of definition is contained in \(F\).\par
\section*{Clarifications and unresolved scope}
\subsection*{U1: non-isomorphic versus non-isogenous tori}
Location. Proposition 3.7(3), p. 20; its local choice on p. 21; the first paragraph of the proof of Lemma 3.9, p. 23; the closure argument in Theorem 3.10, p. 27.\par
Exact concern. The stated property is that the irreducible \(\mathbb Q\)-tori \(T_i\) are pairwise non-isomorphic. The proof on p. 23 then treats \(X^*(T_i)_\mathbb Q\) as pairwise non-isomorphic. In general, rational character modules classify tori up to isogeny, not isomorphism. The full-product conclusion requires non-isogeny, or a further argument ruling out diagonal subgroups.\par
Here is an exact lattice test of the logical distinction. Let the signed permutation group on two coordinates act on \[L_0=\mathbb Z^2,\qquad
 L_1=\{(m,n)\in\mathbb Z^2:m+n\equiv0\pmod2\}.\] Both rationalizations are the same irreducible representation. A basis matrix for \(L_1\) is \[B=\begin{pmatrix}1&1\\1&-1\end{pmatrix},\qquad |\det B|=2.\] The common commutant of the sign change \(\operatorname{diag}(-1,1)\) and coordinate swap in \(M_2(\mathbb Q)\) consists only of scalars. An equivariant lattice isomorphism would therefore be scalar. Equality on each coordinate axis would force its absolute scalar to be \(2\), but \(2\mathbb Z^2\) does not contain \((1,1)\in L_1\). Thus the integral modules are not isomorphic although the rational ones are. A signed-permutation Galois action realizes the same distinction for tori. This calculation tests the inference; it is not a construction of a Shimura datum satisfying every hypothesis of Lemma 3.9.\par
Counterevidence and local repair. The construction on p. 21 can explicitly choose, for every pair \((i,j)\), a prime at which both ambient groups split, and take a split maximal torus on one side and an anisotropic maximal torus on the other. The cited existence input already used there permits these choices. Local split rank is invariant under isogeny; these two tori are therefore not isogenous. The local twisting preserves their types at that prime. Keeping this stronger property throughout gives pairwise non-isomorphic rational character modules, and the semisimple-module argument on pp. 23 and 27 then works.\par
Scope and disposition. This verifies a concrete strengthening of the construction sufficient for the main argument, conditional on the same external existence inputs. It does not prove that the standalone weaker hypotheses of Lemma 3.9 always imply the intermediate assertion printed on p. 23. That question is Needs Evidence; no all-hypothesis counterexample to the lemma or main theorem is claimed. The report does not infer a Major or Minor finding solely from the abstract lattice example.\par
\subsection*{U2: the torus inside a fixed hyperspecial subgroup}
Location. Lemma 3.9, last paragraph of p. 24 and first paragraph of p. 25. The equality \[G_{\mathbb Z_p}(\mathbb Z_p)\cap T_x(\mathbb Q_p)=T_x(\mathbb Q_p)^0\] is used to infer that the special-point field is unramified, and then to take the uniformizer to be \(p\).\par
Exact check. An unramified torus as an abstract local group need not be in good integral position relative to a prescribed reductive model. For \(p=5\), in \(\operatorname{GL}_2(\mathbb Q_5)\) let \(D\) be the diagonal torus, \[u=\begin{pmatrix}1&1/5\\0&1\end{pmatrix},\qquad T=uDu^{-1}.\] This is a split, hence unramified, maximal torus, contained in a rational Borel. Its maximal compact subgroup contains \[u\begin{pmatrix}1&0\\0&2\end{pmatrix}u^{-1}
 =\begin{pmatrix}1&1/5\\0&2\end{pmatrix},\] which is not in \(\operatorname{GL}_2(\mathbb Z_5)\). The printed general inference therefore needs an integral embedding condition. The author’s public 2022 draft of the cited Getz–Hahn book states Lemma A.5.2 for a torus and Borel subgroup scheme of the reductive integral model. That draft is supplementary evidence only; equivalence with every detail of the 2024 edition was not assumed.\par
This example verifies the local distinction, but has not been embedded here in a full higher-rank Shimura datum with all the global irreducibility and level conditions. Accordingly it is not presented as a counterexample to Lemma 3.9’s infinite-order conclusion.\par
Counterevidence and repair of the required valuation step. Unramifiedness of the special-point field is not needed for the infinite-order argument. Retain an arbitrary uniformizer \(\varpi\) of \(L=E(sK_0)_v\), let \(e=e(L/\mathbb Q_p)\), and normalize \(v_p(p)=1\). Since \(\mu_x\) is defined over \(L\), \[v_p\!\left(\chi\!\left(N_{L/\mathbb Q_p}(\mu_x(\varpi))\right)\right)
 =\frac{1}{e}\sum_{\sigma:L\hookrightarrow\overline{\mathbb Q}_p}
       \langle\chi,{}^\sigma\mu_x\rangle.\] All conjugate cocharacters are dominant for the rational Borel. On each nontrivial adjoint factor there is a positive root with positive pairing for at least one summand, and all summands for that root are nonnegative. Thus the displayed valuation is positive. The compact right side of Equation (7) has zero character valuations, so the projection of \(q\) has infinite order. Equation (6) transfers this to the Frobenius element. This reasoning uses compactness, not maximality of that compact subgroup.\par
Scope and disposition. The norm-valuation replacement is checked and preserves the needed infinite-order conclusion once the preceding torus and reciprocity steps hold. The separate assertion that the special-point field is unramified under exactly the printed assumptions remains Needs Evidence. A complete uniform correction of all preparatory claims and all external inputs was not certified. No failure of the main compatibility conclusion was established.\par
\subsection*{A small product-group clarification}
In the first sentence of the proof of Lemma 3.6 (p. 19), an arbitrary open subgroup of a product need not itself decompose as the product of its projections. For example, inverse images of a diagonal subgroup under finite quotient maps give open subgroups with coupled factors. What is sufficient is to choose a product of open subgroups contained in \(K_v\). Such a product exists by the product topology. Applying the factorwise argument to that smaller product gives conjugators in the original \(K_v\). This is a local proof clarification with a complete repair and no change to the lemma’s conclusion.\par
\section*{Disposition and verification limits}
Item | Status\par
Confirmed severity | No confirmed Major or Minor finding.\par
Unresolved scope | U1: weaker torus hypothesis; U2: claimed unramified special-point field and fixed-model embedding.\par
Repair checks | Non-isogeny can be imposed in the construction; the norm-valuation step works with ramification.\par
Main theorems | No all-hypothesis counterexample established; full proof certification not claimed.\par
Adoption | No author adoption requested or assumed.\par
Acceptance | Pending; human reviews empty.\par
All 31 original pages were actually inspected. This is broader page coverage than a claim-focused selection, but does not amount to reproving every external theorem. Outstanding dependencies include the precise companion theorem and its tameness/monodromy consequences; tame specialization; geometric ramification detection; Margulis superrigidity; uniform integrality of special points; torus weak approximation and generic Galois groups; and the global cohomological twisting/Hasse-principle input. The Serre-group discussion in Remark 3.11 was read but not independently reconstructed from the full cited book. The limited title/DOI correction search located no article-specific erratum; that does not prove absence of corrections.\par
\section*{Supporting checks}
It computed the displayed conjugated unit, the scalar commutant, the index-two lattice basis and integral action matrices. The valuation formula and the proofs above were checked analytically; they are not claimed to have been proved by the program. Exact outputs and process receipts are in `evidence/independent\_checks.json` and `evidence/independent\_execution\_log.json`.\par
0691a57a2aae841419ee97b8fac68aa8f3d499ed884933422a388106febab576\par
The reviewing assistant performed the scientific analysis, calculation, compilation and page inspection. The main dispatcher supplied scheduling and attributed webpage UI observation only. Scientific decision remains pending and human reviews remain empty.\par
9 C. Klevdal and S. Patrikis, Compatibility of canonical \(\ell\)-adic local systems on adjoint Shimura varieties, Invent. Math. 242 (2025), 305–335. DOI 10.1007/s00222-025-01357-6. Bound publisher PDF. J. R. Getz and H. Hahn, An Introduction to Automorphic Representations with a view toward Trace Formulae, author’s draft dated 22 April 2022, Appendix A.5, printed p. 482, Lemmas A.5.1–A.5.2. Public author source: https://sites.duke.edu/jgetz/files/2022/04/Graduate\_Text.pdf. Only the relevant local hypotheses were checked; this draft is not silently identified with the cited 2024 edition.\par
\end{document}
